\chapter{Introdução}

\section{Conceito de comunicação institucional}

A comunicação do CRP~SP está pautada pelos princípios da ética, da responsabilidade pública, da transparência, da credibilidade e da colaboração. Materializa a voz e `persona' institucionais e, por meio dela, a autarquia mantém diálogo constante com a categoria e a sociedade. Empenha-se em comunicar a diversidade e a representatividade (racial, étnica, de gênero, de corpos, etária, de territórios). Seus produtos adotam linguagem inclusiva (simples, gendrada, antirracista, anticapacitista e com recursos de acessibilidade).

\section{Princípios institucionais}

Os princípios institucionais são derivados das normativas referentes ao Sistema Conselhos de Psicologia e, principalmente, do \emph{Código de ética profissional da/o psicóloga/o}.

Além dos princípios elencados a seguir, o Jornal Psi está sujeito ao \emph{Código de ética dos jornalistas brasileiros}.

\subsection{Respeito e promoção dos Direitos Humanos}

\subsubsection{Linguagem inclusiva}

\subsubsection{Acessibilidade}

Para que pessoas com dificuldades sensoriais ou cognitivas possam usufruir de todos os produtos comunicacionais e/ou editoriais do CRP~SP, é necessário planejamento e investimento em acessibilidade.

Nossos produtos são publicados em diversos formatos: no `site', como livro digital, na forma de postagens. Cada um desses formatos tem diretrizes específicas de acessibilidade: algumas delas são apresentadas no anexo B, seção  \ref{acessibilidade}.

\subsection{Efetivação de funções precípuas do Conselho Regional de Psicologia}

\subsection{Princípios constitucionais aplicáveis à comunicação pública}

O \href{https://www.planalto.gov.br/ccivil_03/decreto-lei/del0200.htm}{Decreto-Lei nº~200}, de 25 de fevereiro de 1967, define \textbf{autarquia} como ``o serviço autônomo, criado por lei, com personalidade jurídica, patrimônio e receita próprios, para executar atividades típicas da Administração Pública, que requeiram para seu melhor funcionamento, gestão administrativa e fincanceira descentralizada.'' Como autarquia, portanto, o CRP~SP está subordinado às normas e princípios que regem a administração pública.

Por sua vez, o artigo 37 da \href{https://www.planalto.gov.br/ccivil_03/constituicao/constituicao.htm}{Constituição Federal} prevê que ``A administração pública direta e indireta de qualquer dos Poderes da União, dos Estados, do Distrito Federal e dos Municípios obedecerá os princípios de legalidade, impessoalidade, moralidade, publicidade e eficiência \xelip''. A aplicação desses princípios à comunicação do CRP~SP é elucidada a seguir.

\subsubsection{Legalidade}

Diferentemente de instituições privadas, às quais é permitido fazer tudo aquilo que não seja expressamente proibido em lei, às autarquias é permitido fazer \textbf{apenas} o que a lei prevê. De acordo com Paulo e Alexandrino,

\begin{quote}
\xelip na prática de um ato individual, o agente público está obrigado a observar não só a lei e os princípios jurídicos, mas também os decretos, as portarias, as instruções normativas, os pareceres normativos, em suma, os atos administrativos gerais que sejam pertinentes àquela situação concreta com que \xelip se depara.\footnote{PAULO, Vicente; ALEXANDRINO, Marcelo. \textbf{Direito constitucional descomplicado}. 8.~ed. Rio de Janeiro: Forense; São Paulo: Método, 2012.}
\end{quote}

A obediência ao princípio da legalidade deve, ainda, levar em conta a hierarquia do ordenamento jurídico brasileiro, que tem a Constituição Federal como lei maior, e abaixo dela as leis ordinárias --- das quais destacamos a \href{https://transparencia.cfp.org.br/wp-content/uploads/2008/08/lei_1971_57661.pdf}{Lei nº~5.766}, de 20 de dezembro de 1971, que cria o Conselho Federal e os Conselhos Regionais de Psicologia, e a \href{https://transparencia.cfp.org.br/wp-content/uploads/2008/08/lei_1998_96491.pdf}{Lei nº~9.649}, de 27 de maio de 1998, que dispõe, entre outras coisas, sobre os serviços de fiscalização de profissões regulamentadas --- e as normativas próprias da autarquia (portarias, resoluções etc.).

\subsubsection{Impessoalidade}

O princípio da impessoalidade determina que a atuação dos órgãos da administração pública tem por finalidade atender ao interesse público. No caso da comunicação do CRP~SP, isso implica tanto o atendimento das demandas de profissionais de Psicologia quanto a defesa da categoria perante a sociedade. A comunicação institucional da autarquia deve, portanto, ser útil às psicólogas e aos psicólogos e responsável em sua relação com a população em geral, pois representa a face oficial da Psicologia no estado de São Paulo.

Como consequência do compromisso assumido com o interesse público, é vedada a utilização da `persona' institucional para a promoção de indivíduos ou grupos. A comunicação institucional não pode fugir ao princípio da legalidade, devendo restringir-se às situações previstas nas normativas da instituição e na legislação pátria; segundo o § 1º do artigo 37 da Constituição Federal,

\begin{quote}
	A publicidade dos atos, programas, obras, serviços e campanhas dos órgãos públicos deverá ter caráter educativo, informativo ou de orientação social, dela não podendo constar nomes, símbolos ou imagens que caracterizem promoção pessoal de autoridades ou servidores públicos.\footnote{BRASIL. [Constituição (1988)]. \textbf{Constituição da República Federativa do Brasil}. Brasília, DF: Presidência da República, [2024].}
\end{quote}

\subsubsection{Moralidade}

\subsubsection{Publicidade}

\subsubsection{Eficiência}


\section{Implementação e gestão desta Política}

Esta Política deve ser complementada por um \emph{Manual de comunicação}, a ser elaborado por agência para isso capacitada, contendo as seções a seguir.

\subsection{Seções gerais}

Composta de guias a serem utilizados por todas as instâncias e unidades do CRP~SP, inclusive subsedes.

\subsubsection{Redação}

\begin{itemize}
	\item
	Linguagem simples;
	\item
	linguagem inclusiva;
	\item
	tom e extensão dos textos para diferentes mídias;
	\item
	estrutura textual para diferentes produtos;
\end{itemize}

\subsubsection{Audiovisual}

\begin{itemize}
	\item
	Retratos e fotografia de eventos;
\end{itemize}

\subsection{Editoração}

\subsubsection{Revisão textual e de linguagem}

\paragraph{Padrões editoriais}

\begin{itemize}
	\item
	Normatização de textos de acordo com a ABNT,
	\item
	\gls{padronizacao} institucional;
	\item
	boas práticas de acessibilidade para produtos digitais e/ou impressos;
\end{itemize}

\subsubsection{Diagramação}

\subsubsection{Catalogação (ISBN)}

\subsubsection{Preparação de produtos editoriais}

\begin{itemize}
	\item
	Manuais;
	\item
	guias;
	\item
	notas técnicas;
	\item
	cadernos temáticos;
	\item
	relatórios institucionais;
	\item
	anais;
	\item
	cartilhas.
\end{itemize}

\subsection{Produção gráfica}

\begin{itemize}
	\item
	Panfletos;
	\item
	camisetas;
	\item
	adesivos;
	\item
	faixas;
	\item
	`banners';
	\item
	testeiras;
	\item
	crachás;
	\item
	sacolas;
	\item
	placas/sinalização;
	\item
	selos.
\end{itemize}

\subsection{Marketing}

\subsubsection{Planos de comunicação integrada}

\subsubsection{Identidade visual}

A \gls{idv} do \ac{crp}

\subsection{Comunicação externa}

\subsubsection{Relacionamento com a imprensa}

\subsubsection{Mídias sociais}

\subsection{Eventos}

\subsubsection{Eventos nos territórios}
